\section{Set presentations, depth bounds, and a box-label obstruction}
\label{sec:set-presentations}

The representation question has two different parts. Changing arbitrary finite
vertex sets into tagged internal and boundary vertices is now proved harmless.
So is recursively making that change inside binder-annotated boxes, and
normalizing a finite family of different depth allowances to one common bound.
However, these results do not yet justify forgetting the binder annotations
from box labels. The final subsection gives a checked counterexample to that
particular proposed forgetful translation. Section~\ref{sec:paper-encoding} supplies the freshening translation that
resolves this obstruction and completes the representation obligation.

\subsection{Arbitrary vertex sets and ordered boundaries}
\termdef{set-presentation}
\begin{definition}[Set presentation]
For a signature $\sigma$ and ordered boundary types $I,O$, a set presentation
has finite vertex and edge carriers $V,E$, vertex and edge labels, source and
target lists, and injective boundary maps
$i:\{0,\ldots,|I|-1\}\to V$ and $o:\{0,\ldots,|O|-1\}\to V$.
The boundary images are disjoint and have the prescribed vertex labels.
Source and target lists are separately duplicate-free, and no vertex occurs
in the source lists of two different edges, or in the target lists of two
different edges. A vertex has a source incidence exactly when it is not an
output, and a target incidence exactly when it is not an input. Mapping vertex
labels over an incidence list gives the signature of its edge label.
\end{definition}
This is the fixed-signature form of the paper's outer graph definition.
The two boundary functions encode ordered lists by their indices. Disjointness
and coverage exclude isolated vertices, but allow cycles and a vertex shared
by the source and target of the same edge. The signature formulation implies
the paper's equal-label/equal-signature condition. This implication concerns
the chosen label alphabet; changing that alphabet requires a separate check.
\lean{Graph/SetPresentation}{11}{37}{Graph.SetNet}

\termdef{boundary-normalization}
\begin{definition}[Boundary normalization]
The interior $V^{\circ}$ consists of vertices outside both boundary images.
The normalized vertex carrier is the disjoint union of $V^{\circ}$, input
indices, and output indices. The map $\mathsf{decode}$ sends an interior to
its old vertex and a boundary index to the corresponding boundary vertex.
The inverse candidate $\mathsf{encode}$ first tests input membership, then
output membership, and otherwise records the vertex as interior. Classical
choice supplies the index when membership is known.
\end{definition}
\lean{Graph/SetPresentation}{40}{43}{Graph.SetNet.Interior}

\lean{Graph/SetPresentation}{44}{49}{Graph.SetNet.decode}

\lean{Graph/SetPresentation}{50}{54}{Graph.SetNet.encode}

\begin{lemma}[Inverse maps]
Decoding an encoded vertex returns the vertex. Decoding is injective, and
encoding a decoded normalized vertex returns that normalized vertex.
Consequently the two maps form a bijection.
\end{lemma}
\begin{proof}
The first identity follows by the three membership cases and the specification
of the chosen index. For injectivity, compare the constructors of both
normalized vertices. Two interiors are equal by subtype extensionality;
two inputs or two outputs are equal by boundary injectivity. Mixed cases
contradict either an interior's defining exclusions or boundary disjointness.
Apply this injectivity to the first identity to obtain the other inverse law.
\end{proof}
\lean{Graph/SetPresentation}{56}{64}{Graph.SetNet.decode_encode}

\lean{Graph/SetPresentation}{65}{79}{Graph.SetNet.decode_injective}

\lean{Graph/SetPresentation}{80}{83}{Graph.SetNet.encode_decode}

\lean{Graph/SetPresentation}{84}{85}{Graph.SetNet.vertices}

\begin{lemma}[Labels and boundary classification]
Encoding preserves vertex labels. An encoded vertex is tagged as output
exactly when its old vertex is in the output image; the analogous statement
holds for inputs.
\end{lemma}
\begin{proof}
For label preservation first prove the claim for decoding, by the three
constructors and the two boundary typing laws, then use the inverse identity.
For boundary classification, a tag supplies its index; conversely, an old
boundary vertex is the decoding of that tag, so the second inverse identity
forces the correct tag. No assumption of decidable carrier equality is needed.
\end{proof}
\lean{Graph/SetPresentation}{87}{96}{Graph.SetNet.encode_label}

\lean{Graph/SetPresentation}{97}{110}{Graph.SetNet.output_iff}

\lean{Graph/SetPresentation}{111}{123}{Graph.SetNet.input_iff}

\begin{definition}[Normalized graph]
Keep the edge carrier and edge labels, use $V^{\circ}$ as the internal carrier,
and map $\mathsf{encode}$ over every source and target list. Interior finiteness
is inherited by taking a subtype of the old finite vertex enumeration.
\end{definition}
\begin{proof}[Verification of the graph laws]
Injectivity of encoding preserves duplicate-free incidence lists. Two source
occurrences at a normalized vertex decode to the same old vertex, so source
uniqueness identifies their edges; the target argument is identical. Inverse
maps transfer list membership, and the boundary-classification lemma transfers
the coverage equivalences. Finally, mapping labels after encoding is the same
as mapping the old labels, so the old signature equations prove typing.
\end{proof}
\lean{Graph/SetPresentation}{126}{172}{Graph.SetNet.toNet}

\subsection{Outer round trips and transport of isomorphisms}
\begin{definition}[Set-graph isomorphism and exposure]
A set-graph isomorphism comprises bijections of vertices and edges preserving
labels, boundary indices, and the ordered incidence lists. Exposing an existing
\texttt{Net} uses its entire tagged vertex carrier as one ordinary finite
carrier; the input and output maps are its boundary constructors.
\end{definition}
The exposure enumeration concatenates the internal enumeration and the two
finite index enumerations. The constructor cases prove coverage, disjointness,
and boundary typing; all remaining incidence laws are inherited.
\lean{Graph/SetPresentation}{175}{189}{Graph.SetNet.Iso}

\lean{Graph/SetPresentation}{190}{221}{Graph.Net.toSet}

\begin{theorem}[Outer round trips]
Every set presentation is isomorphic to its normalization followed by exposure.
Every tagged graph is isomorphic to its exposure followed by normalization.
\end{theorem}
\begin{proof}
For the first direction use the encoding bijection on vertices and the identity
on edges. Its inverse laws, label preservation, and boundary identities already
supply all fields. For the second direction, the interior of an exposed tagged
graph contains precisely its original internal vertices. The bijection sends
an old internal vertex to its internal tag with the boundary exclusions.
Its inverse eliminates impossible boundary cases. Encoding an arbitrary old
vertex agrees with mapping this interior bijection over its port constructor;
apply that equality to every source and target list to finish the isomorphism.
\end{proof}
\lean{Graph/SetPresentation}{224}{234}{Graph.Net.set_roundtrip}

\lean{Graph/SetPresentation}{235}{248}{Graph.Net.interiorBijection}

\lean{Graph/SetPresentation}{249}{254}{Graph.Net.encode_toSet}

\lean{Graph/SetPresentation}{255}{266}{Graph.Net.roundtrip}

\begin{definition}[Related-label isomorphism and relabelling]
A related-label set isomorphism replaces literal edge-label equality by a
supplied relation $R$; all other conditions remain unchanged. Relabelling
along a signature-preserving function changes only edge labels. The new typing
proof composes the old signature equation with the signature-preservation
identity; all vertex and incidence fields are unchanged.
\end{definition}
\lean{Graph/SetPresentation}{273}{283}{Graph.SetNet.RelIso}

\lean{Graph/SetPresentation}{284}{309}{Graph.SetNet.relabel}

\begin{lemma}[Normalization of related-label isomorphisms]
A related-label set isomorphism induces a related-label isomorphism between
the normalized graphs, with the same edge bijection.
\end{lemma}
\begin{proof}
Boundary membership is preserved by the boundary equations and reflected by
vertex injectivity. Hence the vertex bijection restricts to the interiors;
its inverse uses the inverse vertex map and the same boundary equations.
Encoding commutes with that restricted bijection: decode both sides, inspect
the port constructor, and use the inverse identities. Mapping this commuting
square over each incidence list proves the two required list equations.
Vertex labels and related edge labels come directly from the original
isomorphism. These are the declarations \texttt{input\_iff},
\texttt{output\_iff}, \texttt{interiors}, \texttt{encode\_map}, and
\texttt{RelIso.toNet}, respectively.
\end{proof}
\lean{Graph/SetPresentation}{318}{323}{Graph.SetNet.RelIso.input_iff}

\lean{Graph/SetPresentation}{324}{330}{Graph.SetNet.RelIso.output_iff}

\lean{Graph/SetPresentation}{331}{340}{Graph.SetNet.RelIso.interiors}

\lean{Graph/SetPresentation}{341}{353}{Graph.SetNet.RelIso.encode_map}

\lean{Graph/SetPresentation}{354}{371}{Graph.SetNet.RelIso.toNet}

\subsection{Recursive annotated boxes}
\termdef{annotated-set-hierarchy}
\begin{definition}[Annotated set hierarchy]
An annotated set box records a bound count $b$, a free count $f$, and a set
graph body with interface $\star\Rightarrow\star^{b+f}$.
Its external signature is $T^b(\star)\Rightarrow\star^f$.
The depth-zero alphabet contains primitives; the next alphabet adds boxes over
the preceding alphabet. A hypernet is a set presentation over this alphabet.
This intentionally retains the annotations of the existing Lean hierarchy.
\end{definition}
\lean{Graph/SetHierarchy}{10}{15}{Linear.SetHierarchy.SetBox}

\lean{Graph/SetHierarchy}{16}{23}{Linear.SetHierarchy.nextAlphabet}

\lean{Graph/SetHierarchy}{24}{27}{Linear.SetHierarchy.alphabet}

\lean{Graph/SetHierarchy}{28}{28}{Linear.SetHierarchy.Hypernet}

\begin{definition}[Recursive conversion]
Normalization at depth zero uses outer normalization. At a successor depth it
also relabels every box by the same counts and the recursively normalized body.
Exposure performs the inverse-direction construction: expose the outer carrier
and recursively expose every box body. Primitives are unchanged, and the
external signatures agree because neither operation changes the counts.
\end{definition}
\lean{Graph/SetHierarchy}{31}{39}{Linear.SetHierarchy.normalize}

\lean{Graph/SetHierarchy}{40}{46}{Linear.SetHierarchy.expose}

\begin{theorem}[Recursive typed round trip]
Every typed hypernet is recursively isomorphic to its exposure followed by
normalization.
\end{theorem}
\begin{proof}
Induct on the depth. The base is the outer round trip. At a successor, use its
interior bijection, identity edge bijection, and incidence equations. A primitive
label remains literally equal. A box retains both counts, and its bodies are
recursively isomorphic by induction. This supplies exactly the recursive
label relation required by \texttt{Linear.DeepIso}.
\end{proof}
\lean{Graph/SetHierarchy}{49}{65}{Linear.SetHierarchy.typed_roundtrip}

\begin{definition}[Recursive set isomorphism]
Related set boxes have equal binder counts and related bodies at their common
interface. Set-label equivalence is equality for primitives and existence of a
recursive set-graph isomorphism for boxes. A recursive set isomorphism is an
outer related-label isomorphism using that relation.
\end{definition}
\lean{Graph/SetHierarchy}{68}{73}{Linear.SetHierarchy.BoxRelated}

\lean{Graph/SetHierarchy}{74}{80}{Linear.SetHierarchy.LabelEquiv}

\lean{Graph/SetHierarchy}{81}{81}{Linear.SetHierarchy.DeepIso}

\begin{theorem}[Recursive set round trip]
Every annotated set hypernet is recursively isomorphic to its normalization
followed by exposure.
\end{theorem}
\begin{proof}
Again induct on depth. The outer encoding bijection preserves labels and ordered
boundaries. At a box label the counts are unchanged and the induction hypothesis
supplies the body isomorphism. Thus independent vertex renaming at every
nesting level is accounted for, rather than just renaming the outer graph.
\end{proof}
\lean{Graph/SetHierarchy}{85}{112}{Linear.SetHierarchy.set_roundtrip}

\begin{lemma}[Metadata determined by external type]
If two annotated boxes have equal external input and output types, their bound
and free counts agree.
\end{lemma}
\begin{proof}
Equality of the singleton input lists gives equality of the thunk labels, and
injectivity of the thunk constructor gives the bound count. Taking lengths of
the replicated output lists gives the free count. Notice the premise: the
\emph{external type} is fixed. Equality of contents alone does not supply it.
\end{proof}
\lean{Graph/SetHierarchy}{116}{120}{Linear.SetHierarchy.metadata_unique}

\begin{theorem}[Transport and reflection of recursive isomorphism]
Normalization transports recursive set isomorphisms; exposure transports
recursive typed isomorphisms. Two typed hypernets are recursively isomorphic
if and only if their exposures are recursively set-isomorphic.
\end{theorem}
\begin{proof}
For normalization, restrict the outer vertex bijection to interiors as above.
Induct on depth for edge labels: primitive pairs retain their equality,
primitive/box pairs are impossible, and a related box pair supplies the body
isomorphism to which induction applies. Exposure uses the port bijection of
the original internal-vertex bijection and the same recursive argument for
labels. For reflection, normalize the proposed set isomorphism and compose
with the typed round trip on the left and its inverse on the right. The forward
direction is exposure transport itself.
\end{proof}
\lean{Graph/SetHierarchy}{145}{175}{Linear.SetHierarchy.normalize_iso}

\lean{Graph/SetHierarchy}{176}{223}{Linear.SetHierarchy.expose_iso}

\lean{Graph/SetHierarchy}{224}{229}{Linear.SetHierarchy.expose_iso_iff}

\subsection{Coherence of depth allowances}
\termdef{depth-promotion}
\begin{definition}[Depth promotion]
Promotion increases the permitted depth by one while preserving the outer
carrier, incidence and primitive labels; each box body is recursively promoted.
Repeated promotion, written $\mathsf{raiseBy}(n,k)$, increases a bound $n$ to
$n+k$. There is a corresponding operation on typed hypernets.
\end{definition}
\lean{Graph/SetHierarchy}{123}{129}{Linear.SetHierarchy.raise}
\lean{Graph/SetHierarchy}{232}{236}{Linear.SetHierarchy.raiseBy}

\lean{Graph/SetHierarchy}{237}{240}{Linear.SetHierarchy.typedRaiseBy}

\begin{theorem}[Normalization commutes with promotion]
Normalizing after one promotion equals promoting after normalization.
The same equality holds for any finite number of promotions.
\end{theorem}
\begin{proof}
Induct on the original depth for one promotion. At depth zero the two graph
records are definitionally equal. At a successor all outer data still agree;
primitive labels agree immediately and the box-body fields agree by induction.
Proof fields are irrelevant to equality. Induction on the number of promotions
then proves the iterated equation using the one-step equation at the current
bound. These are equalities of the resulting typed graphs, not merely
existence of some isomorphism.
\end{proof}
\lean{Graph/SetHierarchy}{132}{142}{Linear.SetHierarchy.normalize_raise}
\lean{Graph/SetHierarchy}{243}{249}{Linear.SetHierarchy.normalize_raiseBy}

\termdef{uniform-box-family}
\begin{definition}[Uniformization of a finite box family]
A box family indexed by a finite carrier $E$ assigns each member its own depth,
bound count, free count, and body. Take the maximum of the depths in an
exhaustive finite enumeration of $E$, with maximum zero for an empty list.
Pad each member from its original depth to that common maximum.
\end{definition}
\lean{Graph/SetHierarchy}{252}{258}{Linear.SetHierarchy.BoxFamily}

\lean{Graph/SetHierarchy}{259}{260}{Linear.SetHierarchy.maximum}
\lean{Graph/SetHierarchy}{272}{273}{Linear.SetHierarchy.BoxFamily.common}
\lean{Graph/SetHierarchy}{280}{283}{Linear.SetHierarchy.BoxFamily.uniform}

\begin{lemma}[Common bound and conversion]
Each member's depth is at most the common maximum. Normalizing its uniformized
body equals promoting its normalized original body to that same maximum.
\end{lemma}
\begin{proof}
Membership in a list implies that an entry is bounded by its recursively
computed maximum: for the head use the left maximum inequality, and for a tail
member combine induction with the right maximum inequality. Exhaustiveness
puts every family member in the mapped enumeration, even if the enumeration
contains duplicates. Hence $n+(N-n)=N$, which supplies the interface-preserving
depth transport used in uniformization. Normalization commutes with this
transport by equality elimination. Apply the iterated promotion equation to
finish the claimed equality.
\end{proof}
\lean{Graph/SetHierarchy}{263}{269}{Linear.SetHierarchy.le_maximum}
\lean{Graph/SetHierarchy}{275}{277}{Linear.SetHierarchy.BoxFamily.le_common}
\lean{Graph/SetHierarchy}{287}{297}{Linear.SetHierarchy.BoxFamily.uniform_normalize}
This proves that a finite family of separately bounded annotated contents can
be represented at a common bound. It does not yet formalize a translation of
an undecorated graph-label alphabet from the paper; the distinction below is
independent of finite-depth padding.

\subsection{The obstruction to naive binder erasure}
\termdef{body-consistency}
\begin{definition}[Body consistency for binder erasure]
A successor-depth typed graph is body-consistent when any two box edges whose
contents are literally equal also have equal bound and free counts. The Lean
premise uses heterogeneous equality so contents at definitionally different
interface expressions can be compared. This is the condition needed for the
naive operation that forgets binder annotations but leaves contents unchanged
to respect the paper's equal-label/equal-signature requirement.
\end{definition}
\lean{Graph/SetHierarchyExamples}{11}{15}{Linear.SetHierarchyExamples.BodyConsistent}

\begin{proposition}[Checked obstruction to naive binder erasure]
There is a well-formed typed graph which is not body-consistent.
\end{proposition}
\begin{proof}
Let $V:\star\Rightarrow\star$ be the unary variable corolla. Use this exact
body in two boxes: one with $(b,f)=(0,1)$, and one with $(b,f)=(1,0)$.
Their body records agree, but their external signatures are respectively
$T^0(\star)\Rightarrow\star$ and $T^1(\star)\Rightarrow\epsilon$.
Construct a graph with two input vertices of these thunk types and one term
output. The first edge connects the first input to the output; the second
edge has only the second input and no targets. The graph has no internal
vertices. Direct finite case analysis verifies all incidence, coverage and
typing conditions. Applying body consistency to its two edges would assert
$0=1$, a contradiction.
\end{proof}
\lean{Graph/SetHierarchyExamples}{18}{18}{Linear.SetHierarchyExamples.exposed}

\lean{Graph/SetHierarchyExamples}{19}{20}{Linear.SetHierarchyExamples.bound}

\lean{Graph/SetHierarchyExamples}{21}{22}{Linear.SetHierarchyExamples.identical_contents}

\lean{Graph/SetHierarchyExamples}{23}{25}{Linear.SetHierarchyExamples.different_signatures}
\lean{Graph/SetHierarchyExamples}{28}{74}{Linear.SetHierarchyExamples.twoBoxes}
\lean{Graph/SetHierarchyExamples}{77}{80}{Linear.SetHierarchyExamples.forgetting_binders_fails}

The annotated normalization and round-trip theorems still apply to this
counterexample, as checked by \texttt{twoBoxes\_roundtrip}. A second example
checks a box with two holes, exercising recursive contents rather than just
outer incidence. Finally, \texttt{promoted\_box} checks normalization of that
box after two promotions. These examples test the exact scope of the theorems;
none silently assumes body consistency.
\lean{Graph/SetHierarchyExamples}{84}{88}{Linear.SetHierarchyExamples.twoBoxes_roundtrip}

\lean{Graph/SetHierarchyExamples}{89}{93}{Linear.SetHierarchyExamples.nested_roundtrip}

\lean{Graph/SetHierarchyExamples}{94}{100}{Linear.SetHierarchyExamples.promoted_box}

\paragraph{Relation to the completed representation theorem.}
The results above prove arbitrary-vertex boundary normalization, recursive
round trips for the annotated hierarchy, isomorphism transport, and finite
mixed-depth uniformization. The counterexample refutes only naive binder erasure.
Section~\ref{sec:paper-encoding} now gives the required freshening translation:
it changes actual vertex names inside contents, proves both recursive round
trips and quotient inverse laws, and satisfies the literal graph-label signature
condition. No body-consistency restriction is imposed on the typed graphs.
